\documentclass[11pt]{article}
\usepackage[margin=2cm]{geometry}
\usepackage{amsmath,amssymb}
\usepackage{booktabs,array,longtable}
\usepackage{listings}
\usepackage[dvipsnames]{xcolor}
\usepackage{hyperref}
\hypersetup{hidelinks,colorlinks,linkcolor=NavyBlue,citecolor=NavyBlue}
\usepackage{fancyhdr}
\usepackage[T1]{fontenc}

\definecolor{shellbg}{RGB}{248,248,248}
\lstset{
  backgroundcolor=\color{shellbg},
  basicstyle=\ttfamily\scriptsize,
  breaklines=true,
  frame=single, framerule=0.3pt, rulecolor=\color{gray!40},
  xleftmargin=4pt, xrightmargin=4pt
}

\pagestyle{fancy}\fancyhf{}
\lhead{\small\textit{Browndye2 vs PySTARC --- Source-Level Comparison}}
\rhead{\small\textit{\thepage}}
\renewcommand{\headrulewidth}{0.4pt}

\newcommand{\kBT}{k_{\mathrm{B}}T}
\newcommand{\BD}{\textsc{BD2}}
\newcommand{\PS}{\textsc{PySTARC}}
\newcommand{\same}{\textcolor{ForestGreen}{\textbf{SAME}}}
\newcommand{\improved}{\textcolor{Blue}{\textbf{IMPROVED}}}
\newcommand{\changed}{\textcolor{Orange}{\textbf{CHANGED}}}
\newcommand{\notimpl}{\textcolor{Red}{\textbf{NOT IMPL}}}

\title{\Large Browndye2 vs PySTARC\\[0.3em]
\large Comprehensive Source-Level Physics Comparison}
\author{PySTARC Development --- Internal Audit Document}
\date{\today}

\begin{document}
\maketitle
\tableofcontents
\newpage

%% ═══════════════════════════════════════════════════════════════
\section{Overview}

This document provides a line-by-line comparison of every physics
component in Browndye2 (BD2, Gary Huber, C++/OCaml) and PySTARC
(Python/CuPy). For each component, we identify the BD2 source file,
the PySTARC equivalent, the equations used, and whether the physics
is identical, improved, or changed.

\subsection{Verdict legend}
\begin{itemize}
\item \same\ --- identical physics, same equations, verified numerically
\item \improved\ --- same physics intent, better implementation
\item \changed\ --- different approach, justified by equivalence argument
\item \notimpl\ --- BD2 feature not yet in PySTARC
\end{itemize}


%% ═══════════════════════════════════════════════════════════════
\section{Physical constants}

\begin{tabular}{lll}
\toprule
Constant & BD2 & PySTARC \\
\midrule
Water viscosity & 0.243 (\texttt{physical\_constants.hh}) & 0.243 \\
 & (units: internal, = 0.01 poise) & (same internal units) \\
Vacuum permittivity $\varepsilon_0$ & 0.000142 $e^2/(\kBT\cdot\text{\AA})$ & 0.000142 \\
$4\pi$ & \texttt{pi4} in \texttt{pi.hh} & \texttt{4*np.pi} \\
$6\pi$ & \texttt{pi6} in \texttt{pi.hh} & \texttt{6*np.pi} \\
$\kBT$ at 298.15\,K & 1.0 (implicit) & 1.0 (explicit scaling for $T \neq 298.15$) \\
\bottomrule
\end{tabular}

\textbf{Verdict:} \same. All constants match exactly.


%% ═══════════════════════════════════════════════════════════════
\section{Electrostatic grid computation (APBS)}

\subsection{BD2: \texttt{src/aux/make\_apbs\_inputs.cc}, \texttt{run\_apbs\_inputs.cc}}
BD2 generates APBS input files for two nested grids per molecule using
\texttt{mg-manual} with \texttt{chgm spl2} (cubic B-spline charge mapping).
\begin{itemize}
\item Coarse grid: covers Debye length, \texttt{bcfl sdh} (single Debye--H\"uckel)
\item Fine grid: molecular surface, \texttt{bcfl map} (boundary from coarse)
\item Both electrostatic (full LPBE with ions) and Born (vacuum, no ions)
\end{itemize}

\subsection{PySTARC: \texttt{pystarc/pipeline/run\_apbs.py}}
Identical APBS setup:
\begin{itemize}
\item Same \texttt{mg-manual} method, same \texttt{chgm spl2}
\item Same two-level grid hierarchy (coarse for BC, fine for forces)
\item Same electrostatic + Born split
\item \textbf{Addition:} Per-molecule grid sizing (ligand gets tight auto-sized grid; receptor gets override via XML)
\item \textbf{Addition:} Per-level \texttt{dime} support (\texttt{apbs\_coarse\_dime}, \texttt{apbs\_fine\_dime})
\end{itemize}

\textbf{Verdict:} \same\ for APBS physics. \improved\ for grid sizing flexibility.


%% ═══════════════════════════════════════════════════════════════
\section{Grid force interpolation}

\subsection{BD2: \texttt{src/forces/electrostatic/single\_grid.hh} (767 lines)}
Trilinear interpolation of grid values. Gradient computed by central
differences at half-grid-spacing. For atoms outside a grid, falls
through to the next coarser grid or to the Chebyshev blob expansion.

Key code pattern:
\begin{lstlisting}
gradient_of_cube(cube, ax, ay, az)  // trilinear gradient
// ax,ay,az are fractional positions within the cell
\end{lstlisting}

\subsection{PySTARC: \texttt{pystarc/forces/electrostatic/grid\_force.py} (310 lines)}
Same trilinear interpolation with central-difference gradients.
Vectorised over all atoms via NumPy/CuPy broadcasting.

\textbf{Verdict:} \same\ interpolation scheme.
\changed\ for atoms outside grid: PySTARC uses Yukawa multipole (Section~\ref{sec:farfield})
instead of falling through to coarse grid.


%% ═══════════════════════════════════════════════════════════════
\section{Electrostatic force evaluation}

\subsection{BD2: \texttt{src/forces/electrostatic/charge\_field.cc} (153 lines)}
For each ligand atom $i$ at position $\mathbf{r}_i$:
\begin{equation}
\mathbf{F}_i = -q_i \nabla\phi_\mathrm{rec}(\mathbf{r}_i)
\end{equation}
Gradient from finest available grid. If outside fine grid, uses coarse grid.
If outside all grids, uses Chebyshev blob multipole.

\subsection{PySTARC: \texttt{pystarc/forces/gpu\_batch\_engine.py} (627 lines)}
Same per-atom force equation. For atoms outside the fine grid (with 3-spacing
safety margin), uses Yukawa multipole fallback.

\textbf{Key difference:} BD2 uses a two-level grid cascade (fine $\to$ coarse $\to$ Chebyshev).
PySTARC uses fine grid only + Yukawa far-field (coarse grid dropped --- it exists only for APBS boundary conditions).

\textbf{Verdict:} \same\ force equation. \changed\ far-field fallback
(Yukawa multipole vs Chebyshev blobs).


%% ═══════════════════════════════════════════════════════════════
\section{Screened Coulomb (Yukawa) force}

\subsection{BD2: \texttt{src/forces/electrostatic/screened\_coulomb.hh}}
Monopole-only screened Coulomb between two point charges:
\begin{equation}
\mathbf{F} = \frac{q_0 q_1}{4\pi\varepsilon_s}
\frac{(r/\lambda + 1)\,e^{-r/\lambda}}{r^3}\,\mathbf{r}_{01}
\end{equation}
Used in outer propagator for analytical forces between centroids.

\subsection{PySTARC: \texttt{pystarc/forces/multipole\_farfield.py}}
Three-term expansion:
\begin{itemize}
\item Monopole: same as BD2 screened Coulomb (Eq.\ above)
\item Dipole: $\phi_1 \propto (\mathbf{p}\cdot\hat{\mathbf{r}})\,(1+r/\lambda)\,e^{-r/\lambda}/r^2$
\item Quadrupole: $\phi_2 \propto (\hat{\mathbf{r}}^T\mathbf{Q}\hat{\mathbf{r}})\,(1+r/\lambda+r^2/3\lambda^2)\,e^{-r/\lambda}/r^3$
\end{itemize}
Moments precomputed from PQR. For neutral molecules ($Q=0$),
dipole provides the leading non-zero term.

\textbf{Verdict:} \improved. Monopole matches BD2 exactly. Dipole and quadrupole
are additions that improve accuracy for non-spherical charge distributions.


%% ═══════════════════════════════════════════════════════════════
\section{Chebyshev blob expansion}
\label{sec:farfield}

\subsection{BD2: \texttt{src/generic\_algs/cartesian\_multipole.cc} (645 lines)}
Fits effective charges on a sphere to reproduce the potential.
Used when atoms lie outside all APBS grids.

\subsection{PySTARC}
Not implemented. Replaced by the Yukawa multipole expansion (Section above).
For charged systems, monopole dominates at $r > 3a$.
For the two-charged-spheres benchmark, $P_\mathrm{rxn}$ matches
the exact Smoluchowski solution to 0.9\%.

\textbf{Verdict:} \changed. Functionally replaced by Yukawa multipole.


%% ═══════════════════════════════════════════════════════════════
\section{Born desolvation}

\subsection{BD2: \texttt{src/aux/born\_grid.cc}, \texttt{born\_field\_interface.hh}}
Born potential $\phi_\mathrm{born}$ computed by APBS in vacuum
($\varepsilon = 1$, no ions). Force:
\begin{equation}
\mathbf{F}_i^\mathrm{born} = -\alpha\, q_i^2 \nabla\phi_\mathrm{born}(\mathbf{r}_i),
\quad \alpha = \frac{1}{4\pi} \approx 0.0796
\end{equation}
Applied in one direction only (receptor Born at ligand atoms) unless
both-directions is specified in the input.

\subsection{PySTARC: \texttt{pystarc/forces/gpu\_batch\_engine.py}}
Same equation, same $\alpha$. Applied in \textbf{both directions} by default:
\begin{itemize}
\item Direction 1: receptor Born grid, evaluated at ligand atom positions
\item Direction 2: ligand Born grid, evaluated at receptor atom positions
  (with GPU memory chunking for large systems)
\end{itemize}

\textbf{Verdict:} \same\ equation. \improved\ with both-directions default
and GPU memory-safe chunking (500\,MB per batch).


%% ═══════════════════════════════════════════════════════════════
\section{Hydrodynamic radius (Hansen Monte Carlo)}

\subsection{BD2: \texttt{aux/hydro\_params.ml} (OCaml)}
Hansen algorithm: random walkers on the solvent-excluded surface,
mean inverse chord length gives the Stokes radius.
Surface represented as a cubic grid of inside/outside voxels.

\subsection{PySTARC: \texttt{pystarc/hydrodynamics/mc\_hydro\_radius.py} (285 lines)}
Same Hansen algorithm. Surface computed from PQR atom positions + probe radius.
Monte Carlo with configurable number of walkers.
Results cached to \texttt{.r\_hydro\_cache} files.

\textbf{Verdict:} \same\ algorithm. Values match to $<1\%$
(e.g.\ trypsin: BD2 22.50\,\AA, PySTARC 22.50\,\AA).


%% ═══════════════════════════════════════════════════════════════
\section{Diffusion coefficients}

\subsection{BD2: implicit in \texttt{physical\_constants.hh}, \texttt{solvent.cc}}
\begin{align}
D_t &= \frac{\kBT}{6\pi\eta a}, \quad D_r = \frac{3D_t}{4a^2}
\end{align}
$D_0 = D_t^{(1)} + D_t^{(2)}$.

\subsection{PySTARC: \texttt{pystarc/pipeline/geometry.py}}
Same equations. Same viscosity (0.243 internal = 0.01 poise).

\textbf{Verdict:} \same. $D_0$ matches to machine precision.


%% ═══════════════════════════════════════════════════════════════
\section{Romberg integration for $k_b$}

\subsection{BD2: \texttt{aux/romberg.ml} (OCaml)}
Standard Romberg quadrature with Richardson extrapolation.
Integrand: $\exp[V(1/s)/\kBT] / D_\parallel(1/s)$ with change of variables $s = 1/r$.
\begin{equation}
k_b = \frac{4\pi}{\int_0^{1/b} \frac{e^{V(1/s)/\kBT}}{D_\parallel(1/s)}\,ds}
\end{equation}

\subsection{PySTARC: \texttt{pystarc/simulation/outer\_propagator.py} (407 lines)}
Same Romberg quadrature, same change of variables, same integrand.
Uses \texttt{scipy.integrate.romberg} internally.

\textbf{Verdict:} \same. $k_b$ values match:
\begin{itemize}
\item Charged spheres: 57.447 vs 57.447 (exact)
\item Thrombin: 35.733 vs 35.787 ($\Delta = 0.15\%$)
\item Trypsin: 25.632 vs 25.620 ($\Delta = 0.05\%$)
\end{itemize}


%% ═══════════════════════════════════════════════════════════════
\section{Rotne--Prager--Yamakawa hydrodynamic interactions}

\subsection{BD2: \texttt{src/hydrodynamics/rotne\_prager.hh} (322 lines)}
Implements the Zuk et al.\ (2014) RPY tensor, including the overlap regime:
\begin{equation}
D_\parallel(r) = \frac{\kBT}{6\pi\eta}
\left(\frac{1}{a_1}+\frac{1}{a_2}-\frac{3}{r}+\frac{2\bar{a}^2}{r^3}\right)
\end{equation}
for $r > a_1 + a_2$ (non-overlapping). For $r < a_1 + a_2$ (overlap),
uses the Zuk 2014 correction terms.

\subsection{PySTARC: \texttt{pystarc/hydrodynamics/rotne\_prager.py}}
Same Zuk 2014 equations, same overlap corrections.

\textbf{Verdict:} \same. RPY tensor matches exactly.


%% ═══════════════════════════════════════════════════════════════
\section{BD propagator (Ermak--McCammon)}

\subsection{BD2: \texttt{src/motion/do\_bd\_step.cc} (1398 lines)}
\begin{equation}
\mathbf{r}(t+\Delta t) = \mathbf{r}(t) + \frac{D_0}{\kBT}\mathbf{F}\Delta t
+ \sqrt{2D_0\Delta t}\,\mathbf{W}
\end{equation}
$\mathbf{W} \sim \mathcal{N}(\mathbf{0},\mathbf{I})$.
Also handles chain BD (not in PySTARC).

\subsection{PySTARC: \texttt{pystarc/simulation/gpu\_batch\_simulator.py} (1421 lines)}
Same equation, vectorised over all $N$ trajectories simultaneously on GPU.
Single CuPy kernel call updates all active trajectories.

\textbf{Verdict:} \same\ equation. \improved\ with GPU batch execution.


%% ═══════════════════════════════════════════════════════════════
\section{Adaptive time stepping}

\subsection{BD2: \texttt{src/motion/do\_full\_step.cc} (270 lines)}
Three constraints (from \texttt{pair\_dt} function):
\begin{align}
\Delta t_\mathrm{pair} &= \frac{f^2}{2}\frac{r^2}{D}, \quad f = 0.1 \\
\Delta t_\mathrm{force} &\text{ (via Wiener loop backstep)} \\
\Delta t_\mathrm{edge} &\text{ (via outer propagator trigger at } r > 1.1b\text{)}
\end{align}
BD2 uses a growth factor of 1.1$\times$ from the previous $\Delta t$.

\subsection{PySTARC: \texttt{pystarc/simulation/gpu\_batch\_simulator.py}}
Three explicit constraints:
\begin{align}
\Delta t_\mathrm{pair} &= \frac{f^2}{2}\frac{r^2}{D_0}, \quad f = 0.1 \\
\Delta t_\mathrm{force} &= \frac{0.01}{|D_0\mathbf{F}|} \\
\Delta t_\mathrm{edge} &= \frac{\min(r-b, r_\mathrm{esc}-r)^2}{18\,D_0}
\end{align}

\textbf{Verdict:} \same\ pair constraint. \changed\ force and edge constraints
are explicit rather than implicit through the Wiener loop.


%% ═══════════════════════════════════════════════════════════════
\section{Rotational diffusion}

\subsection{BD2: \texttt{src/lib/rotations.cc} (204 lines)}
Lookup table of 181 rotation matrices parameterised by angle.
Interpolates between table entries for the diffusional rotation step.

\subsection{PySTARC: \texttt{pystarc/simulation/diffusional\_rotation.py} (337 lines)}
Exact quaternion algebra. Random rotation axis and angle
$\theta = \sqrt{2D_r\Delta t}\,|\boldsymbol{\xi}|$
composed by quaternion multiplication.

\textbf{Verdict:} \improved. Exact for any rotation angle; no table interpolation error.


%% ═══════════════════════════════════════════════════════════════
\section{Reaction detection}

\subsection{BD2: \texttt{src/generic\_algs/step\_near\_absorbing\_surface.hh} (128 lines),
\texttt{src/generic\_algs/wiener/wiener\_loop.hh} (166 lines)}

Three-stage detection:
\begin{enumerate}
\item Endpoint check: are $n_\mathrm{needed}$ pairs within cutoff?
\item Reaction backstep: if $\sigma = \sqrt{2D\Delta t} > d_\mathrm{rxn}/12$,
  halve $\Delta t$ and retry (Wiener loop).
\item \texttt{step\_near\_absorbing\_surface}: exact 1D Brownian motion near
  an absorbing wall. Uses rejection sampling to generate surviving/absorbed positions.
\end{enumerate}
The Wiener loop can trap when $d_\mathrm{rxn} \to 0$ because $\Delta t_\mathrm{rxn} \to 0$.
BD2 mitigates this with \texttt{minimum\_core\_reaction\_dt} floor.

\subsection{PySTARC: \texttt{pystarc/simulation/gpu\_batch\_simulator.py}}

Two-stage detection:
\begin{enumerate}
\item Endpoint check: same as BD2
\item Brownian bridge: for $x_0, x_1 > 0$ (both endpoints above surface),
  $P_\mathrm{cross} = \exp(-x_0 x_1 / D_\mathrm{eff}\Delta t)$
  with $D_\mathrm{eff} = D_0 + D_r^{(1)}L_1^2 + D_r^{(2)}L_2^2$.
\end{enumerate}
No retry loop, no $\Delta t$ floor needed.

\textbf{Verdict:} \changed. Both catch mid-step crossings.
BD2 uses rejection sampling in 1D; PySTARC uses the exact Brownian
bridge formula. The bridge is $O(1)$ per trajectory and avoids the
$\Delta t \to 0$ trapping problem.


%% ═══════════════════════════════════════════════════════════════
\section{Outer propagator}

\subsection{BD2: \texttt{src/motion/outer\_propagator.cc} (396 lines)}
When $r > q_b \cdot b$ (where $q_b = 1.1$, from \texttt{qb\_factor.hh}):
\begin{enumerate}
\item Run BD with analytical Yukawa forces from $r$ to $q_\mathrm{out} = 20\max(a_1,a_2)$.
\item At $q_\mathrm{out}$: return to $b$ with probability $p_\mathrm{ret} = k_b(b)/k_b(q_\mathrm{out})$,
  or escape.
\item Returned trajectories get a new position on the $b$-sphere and accumulated diffusional rotation.
\end{enumerate}

\subsection{PySTARC: \texttt{pystarc/simulation/gpu\_batch\_simulator.py}}
When $r \ge r_\mathrm{esc} = 2b$:
\begin{enumerate}
\item Apply $p_\mathrm{ret}$ directly (no intermediate BD steps).
\item Yukawa forces are already applied everywhere via the fine grid + multipole fallback.
\item Returned trajectories get random position on $b$-sphere + random rotation.
\end{enumerate}

\textbf{Key difference:} BD2 runs BD steps in the outer zone ($1.1b < r < q_\mathrm{out}$);
PySTARC applies forces everywhere and only uses $p_\mathrm{ret}$ at $r_\mathrm{esc}$.
Both implement the LMZ formulation.

$p_\mathrm{ret}$ formula is identical:
$p_\mathrm{ret} = k_b(b)/k_b(r_\mathrm{esc})$

\textbf{Verdict:} \same\ physics (LMZ). \changed\ implementation
(direct return at escape vs intermediate BD in outer zone).
$k_b$ values match to $<0.2\%$.


%% ═══════════════════════════════════════════════════════════════
\section{Overlap rejection}

\subsection{BD2: \texttt{src/core\_and\_chain/sphere\_overlap.cc}}
Precomputed surface mesh from \texttt{make\_surface\_sphere\_list}.
If any ligand atom enters the receptor's solute dielectric region
(inside the mesh), the step is rejected.

\subsection{PySTARC: \texttt{pystarc/simulation/gpu\_batch\_simulator.py}}
Elastic wall: if $\min_i \|\mathbf{r}_\mathrm{lig} - \mathbf{r}_{\mathrm{rec},i}\| < 1.5$\,\AA,
restore pre-step position and orientation.
GPU-chunked for memory (500\,MB per batch for large receptors).

\textbf{Verdict:} \changed. Both prevent interpenetration.
BD2 uses surface mesh; PySTARC uses distance-based elastic wall.


%% ═══════════════════════════════════════════════════════════════
\section{LJ / WCA steric forces}

\subsection{BD2: \texttt{src/forces/lj\_force\_computer.hh}}
Full Lennard-Jones with ff14SB/GAFF parameters:
$V = \epsilon[(\sigma/r)^{12} - (\sigma/r)^6]$

\subsection{PySTARC: \texttt{pystarc/forces/lj.py} (252 lines)}
WCA (repulsive-only) decomposition:
$V_\mathrm{WCA} = V_\mathrm{LJ} + \epsilon$ for $r < 2^{1/6}\sigma$, else 0.
Uses PQR radii directly ($\sigma_{ij} = R_i + R_j$),
$\epsilon = 0.17\,\kBT$.

\textbf{Verdict:} \changed. Same steric wall position; WCA omits
the attractive well (depth $\sim$0.1\,$\kBT$ per pair --- negligible for $\kon$).


%% ═══════════════════════════════════════════════════════════════
\section{Rate constant computation}

\subsection{BD2: \texttt{aux/compute\_rate\_constant.ml} (OCaml)}
\begin{align}
P_\mathrm{rxn} &= N_\mathrm{reacted} / N_\mathrm{total} \\
\kon &= 6.022 \times 10^8 \times k_b \times P_\mathrm{rxn}
\end{align}
95\% CI: normal approximation $P \pm 2\sqrt{P(1-P)/N}$.

\subsection{PySTARC: \texttt{run\_pystarc.py}, \texttt{pystarc/analysis/convergence.py}}
Same $\kon$ formula. 95\% CI: Wilson score interval:
\begin{equation}
P \in \frac{\hat{p} + z^2/2n \pm z\sqrt{\hat{p}(1-\hat{p})/n + z^2/4n^2}}{1 + z^2/n}
\end{equation}
Valid for any $\hat{p}$ and any $n \ge 1$.

Plus automated convergence analysis: relative SE, cumulative curve,
first-half/second-half split, N-needed estimates.

\textbf{Verdict:} \same\ rate formula. \improved\ confidence interval
(Wilson vs normal approx) and convergence analysis (BD2 has none).


%% ═══════════════════════════════════════════════════════════════
\section{NAM simulator}

\subsection{BD2: \texttt{src/simulation/nam\_simulator.cc} (305 lines)}
Launches $N$ independent trajectories across CPU threads.
Each trajectory: start at $b$-sphere, evolve, check reaction, handle escape/return.
Output: XML file per thread with completion counts.

\subsection{PySTARC: \texttt{pystarc/simulation/gpu\_batch\_simulator.py} (1421 lines)}
All $N$ trajectories stored as GPU arrays. Single loop updates all simultaneously.
Force evaluation batched for large ligands.
14 structured output files (vs BD2's per-thread XML).

\textbf{Verdict:} \same\ NAM algorithm. \improved\ with GPU batch execution
and structured outputs.


%% ═══════════════════════════════════════════════════════════════
\section{Temperature scaling}

\subsection{BD2: implicit}
All energies in units of $\kBT$ at 298.15\,K. No explicit temperature parameter.

\subsection{PySTARC: \texttt{pystarc/pipeline/input\_parser.py}}
Explicit temperature tag. Diffusion and Boltzmann factors scaled by
$T/298.15$. At default $T = 298.15$\,K, scaling = 1.0 (identical to BD2).

\textbf{Verdict:} \same\ at 298.15\,K. \improved\ with explicit temperature support.


%% ═══════════════════════════════════════════════════════════════
\section{Features in BD2 not in PySTARC}

\begin{longtable}{lll}
\toprule
Feature & BD2 source & Status \\
\midrule
Chain BD (flexible) & \texttt{chain/chain\_*.cc} & \notimpl\ (planned) \\
COFFDROP potentials & \texttt{coffdrop.cc} & \notimpl\ (planned) \\
ABSINTH solvation & \texttt{absinth/} & \notimpl\ (planned) \\
Staged reactions & \texttt{staged\_simulator.cc} & \notimpl \\
Weighted ensemble BD & \texttt{we\_simulator.cc} & Partial (in \texttt{we\_simulator.py}) \\
Chebyshev blobs & \texttt{cartesian\_multipole.cc} & Replaced by Yukawa multipole \\
Surface mesh overlap & \texttt{sphere\_overlap.cc} & Replaced by elastic wall \\
Full LJ (attractive) & \texttt{lj\_force\_computer.hh} & Replaced by WCA repulsive \\
Trajectory XML output & \texttt{trajectory\_outputter.cc} & Replaced by 14 structured files \\
\bottomrule
\end{longtable}


%% ═══════════════════════════════════════════════════════════════
\section{Features in PySTARC not in BD2}

\begin{longtable}{ll}
\toprule
Feature & PySTARC source \\
\midrule
GPU batch simulation & \texttt{gpu\_batch\_simulator.py} \\
GPU force evaluation & \texttt{gpu\_batch\_engine.py} \\
Brownian bridge reaction detection & \texttt{gpu\_batch\_simulator.py} \\
Yukawa dipole + quadrupole & \texttt{multipole\_farfield.py} \\
Wilson score CI & \texttt{analysis/convergence.py} \\
Automated convergence analysis & \texttt{analysis/convergence.py} \\
Live progress reporting & \texttt{gpu\_batch\_simulator.py} \\
14 ML-ready output files & \texttt{pipeline/output\_writer.py} \\
Checkpointing + resume & \texttt{gpu\_batch\_simulator.py} \\
Automated APBS pipeline & \texttt{pipeline/run\_apbs.py} \\
GHO atom injection & \texttt{pipeline/gho\_injection.py} \\
XML input parser (40 tags) & \texttt{pipeline/input\_parser.py} \\
WCA steric forces & \texttt{forces/lj.py} \\
Force batch memory management & \texttt{gpu\_batch\_simulator.py} \\
Per-molecule APBS grid sizing & \texttt{pipeline/run\_apbs.py} \\
Temperature scaling & \texttt{pipeline/input\_parser.py} \\
1,031 unit tests & \texttt{tests/} \\
\bottomrule
\end{longtable}


%% ═══════════════════════════════════════════════════════════════
\section{Complete file mapping}

\begin{longtable}{p{5.5cm}p{5.5cm}l}
\toprule
\textbf{BD2 file} & \textbf{PySTARC file} & \textbf{Verdict} \\
\midrule
\endhead
\texttt{physical\_constants.hh} & constants in multiple files & \same \\
\texttt{single\_grid.hh} & \texttt{grid\_force.py} & \same \\
\texttt{charge\_field.cc} & \texttt{gpu\_batch\_engine.py} & \same \\
\texttt{screened\_coulomb.hh} & \texttt{multipole\_farfield.py} & \improved \\
\texttt{born\_grid.cc} & \texttt{gpu\_batch\_engine.py} & \same \\
\texttt{rotne\_prager.hh} & \texttt{rotne\_prager.py} & \same \\
\texttt{hydro\_params.ml} & \texttt{mc\_hydro\_radius.py} & \same \\
\texttt{romberg.ml} & \texttt{outer\_propagator.py} & \same \\
\texttt{do\_bd\_step.cc} & \texttt{gpu\_batch\_simulator.py} & \same \\
\texttt{do\_full\_step.cc} & \texttt{gpu\_batch\_simulator.py} & \same \\
\texttt{outer\_propagator.cc} & \texttt{outer\_propagator.py} & \same \\
\texttt{rotations.cc} & \texttt{diffusional\_rotation.py} & \improved \\
\texttt{wiener\_loop.hh} & (Brownian bridge) & \changed \\
\texttt{step\_near\_absorbing\_surface.hh} & (Brownian bridge) & \changed \\
\texttt{nam\_simulator.cc} & \texttt{gpu\_batch\_simulator.py} & \same \\
\texttt{single\_traj\_simulator.cc} & \texttt{nam\_simulator.py} & \same \\
\texttt{sphere\_overlap.cc} & (elastic wall) & \changed \\
\texttt{lj\_force\_computer.hh} & \texttt{lj.py} & \changed \\
\texttt{cartesian\_multipole.cc} & \texttt{multipole\_farfield.py} & \changed \\
\texttt{compute\_rate\_constant.ml} & \texttt{convergence.py} & \improved \\
\texttt{time\_step\_parameters.cc} & \texttt{input\_parser.py} & \same \\
\texttt{qb\_factor.hh} ($=1.1$) & (escape at $2b$) & \changed \\
\texttt{trajectory\_outputter.cc} & \texttt{output\_writer.py} & \improved \\
\texttt{make\_apbs\_inputs.cc} & \texttt{run\_apbs.py} & \same \\
\texttt{solvent.cc} & \texttt{geometry.py} & \same \\
\bottomrule
\end{longtable}


%% ═══════════════════════════════════════════════════════════════
\section{Numerical validation summary}

\begin{tabular}{llll}
\toprule
Quantity & PySTARC & BD2 / Exact & Match \\
\midrule
$\varepsilon_0$ & 0.000142 & 0.000142 & exact \\
$\eta$ & 0.243 & 0.243 & exact \\
$k_b$ (spheres) & 57.447 & 57.447 & exact \\
$k_b$ (thrombin) & 35.733 & 35.787 & 0.15\% \\
$k_b$ (trypsin) & 25.632 & 25.620 & 0.05\% \\
$P_\mathrm{rxn}$ (spheres) & 0.4462 & 0.4501 (exact) & 0.9\% \\
$D_0$ (thrombin) & 0.01848 & 0.01848 & exact \\
$p_\mathrm{ret}$ (spheres) & 0.5242 & 0.5242 & exact \\
$r_\mathrm{hydro}$ (trypsin rec) & 22.501 & 22.501 & exact \\
$r_\mathrm{hydro}$ (trypsin lig) & 5.056 & 5.056 & exact \\
\bottomrule
\end{tabular}


\end{document}
